%% kapa: exam=1 round=22 year=2011 period=2 form=A subject=accounting subject_ko=회계학 seq=20 number=20
%% source: Q-Net 감정평가사 기출문제 (cst00309) artlSeq=5036790 file=extracted/1008/5036790/2011년 제22회 감정평가사 1차시험 기출문제/2교시.hwp
%% conversion: kordoc 4.24.51 -> tools/gen1.py (auto); review=proofread-gpt-astra; flags=none
\begin{question}{20}
20×6년 1월 1일 (주)서울은 산업합리화 정책의 일환으로 설비자금의 일부를 정부에서 지원받았다. 설비의 취득원가는 ￦500,000이고 정부보조금은 ￦200,000으로 설비취득일에 전액 수령하였다. 동 설비의 내용연수는 5년, 잔존가액은 ￦20,000이며 정액법으로 감가상각한다. (주)서울이 동 설비를 20×9년 1월 1일 ￦150,000에 처분하였을 경우 유형자산처분손익은? (단, 동 설비에 대하여 원가모형을 적용하고 있음, 위의 정부보조금은 상환의무가 없으며 관련자산을 차감하는 방법으로 회계처리함)
\begin{choices}
  \item ￦18,000 이익
  \item ￦18,000 손실
  \item ￦62,000 이익
  \item ￦62,000 손실
  \item ￦92,000 이익
\end{choices}
\end{question}
\answer{1}
